\documentclass{article}
\usepackage{anyfontsize}
\usepackage[UTF8]{ctex} % 如果需要支持中文
\usepackage{fontspec} 
% \setCJKmainfont{SourceHanSansCN-Light}[
%     BoldFont=SourceHanSansCN-Medium
% ]

\usepackage[a4paper, left=0.1in, right=0.1in, top=0.05in, bottom=0.05in]{geometry} % 设置四边页边距为0.1英寸{geometry} % 设置页边距为0.5英寸
\usepackage{enumitem} % 用于自定义列表
\usepackage{amsmath}  % 数学公式支持
\usepackage{tikz}
\usetikzlibrary{arrows.meta, positioning}
\setlength{\parindent}{0pt} % 不缩进
\usepackage{etoolbox} % 用于列表操作
%调整类代码
\usepackage{comment}
\usepackage{tocloft} % 用于定制目录 样式
\usepackage{hyperref} %不需要目录盒子注释这
\usepackage{ifthen}
\usepackage{tikz}
\usetikzlibrary{arrows.meta}
\usepackage{titletoc}
\usepackage{graphicx}
\usepackage{float}

\usepackage{amssymb}  % 提供数学符号，包括\therefore
%目录设定
%快捷键 CMD + / 注释
%  \renewcommand{\cftdot}{} % 隐藏目录中的页码
%  \cftpagenumbersoff{section}
%  \cftpagenumbersoff{subsection}
%  \makeatletter
%  \renewcommand{\@pnumwidth}{0em}  % 设置页码宽度为0
%  \renewcommand{\@tocrmarg}{0em}   % 设置右边距为0
%  \makeatother
 

\usepackage{environ}

\newif\ifshowcontent  % 定义一个条件判断变量
\showcontenttrue    % 设置为 false，隐藏所有 question 环境的内容

\newcounter{question} % 题目计数器
\setcounter{question}{0}
\NewEnviron{question}{%
    \ifshowcontent
        \par\stepcounter{question}\textbf{\thequestion.} \BODY\par % 如果显示内容，则输出
    \fi
}

\NewEnviron{choices}{%
    \ifshowcontent
        \begin{enumerate}[label=\Alph*., leftmargin=*]
            \BODY
        \end{enumerate}\par
    \fi
}

\NewEnviron{correctanswer}{%
    \ifshowcontent
        \textbf{答案:}\BODY\par
    \fi
}



\newenvironment{explanation}
{\textbf{解析:}\par}
{\par}



% 新增考察方面命令
\newcommand{\checklist}[1]{
    \textbf{考察:} #1 \par % 在题目下方显示考察方面
    \addcontentsline{toc}{subsection}{题号\thequestion 考察: #1} % 将考察内容添加到目录
    \vspace{2em} % 添加1em的垂直空间
}

\graphicspath{{images/}} %统一


\begin{document}
 %\fontsize{22}{31}\selectfont
 \tableofcontents % 添加目录
\newpage
%\pagestyle{empty}




\section{考点1:交易账簿的基本面回顾}
\setcounter{question}{0}

\begin{question}
    Which of the following risks is specifically recognized by the incremental risk charge (IRC)?\\
    以下哪一项风险被增量风险收费（IRC）特别识别？
    \end{question}
    
    \begin{choices}
        \item Expected shortfall risk, because it is important to understand the amount of loss potential in the tail\\
        预期短缺风险，因为它很重要，可以理解尾部损失的潜在金额
        \item Jump-to-default risk, as measured by 99.9\% VaR, because a default could cause a significant loss for the bank\\
        跳跃到违约风险，通过99.9\%置信水平的VaR衡量，因为违约可能导致银行重大损失
        \item Equity price risk, because a change in market prices could materially impact mark-to-market accounting for risk\\
        股权价格风险，因为市场价格的变化可能对风险的市场价值产生重大影响
        \item Interest rate risk, as measured by 97.5\% expected shortfall, because an increase in interest rates could cause a significant loss for the bank\\
        利率风险，通过97.5\%置信水平的预期短缺衡量，因为利率上升可能导致银行重大损失
    \end{choices}
    
    \begin{correctanswer}
        B
    \end{correctanswer}
    
    \begin{explanation}
        \textbf{A选项错误。}
        \begin{itemize}
            \item 预期短缺（Expected Shortfall）风险是FRTB（Fundamental Review of the Trading Book）中使用的风险度量方法，而不是IRC的范畴。
        \end{itemize}
        
        \textbf{B选项正确。}
        \begin{itemize}
            \item IRC（增量风险费用）主要识别两种风险：信用利差风险和违约风险。违约风险使用99.9\%置信水平的VaR进行衡量，因为违约事件可能导致银行重大损失。
        \end{itemize}
        
        \textbf{C选项错误。}
        \begin{itemize}
            \item 股价风险属于常规市场风险类别，通过delta/vega/曲率方法或ES进行衡量，不在IRC的考虑范围内。
        \end{itemize}
        
        \textbf{D选项错误。}
        \begin{itemize}
            \item 利率风险也属于常规市场风险类别，使用97.5\%置信水平的ES进行衡量，不在IRC的考虑范围内。
        \end{itemize}
    \end{explanation}
    
    \checklist{IRC识别的两种风险}



    \begin{question}
        Which of the following statements represents a criteria for classifying an asset into the trading book?
        \begin{itemize}
            \item I.The bank must be able to physically trade the asset
            \item II.The risk of the asset must be managed by the bank's trading desk
        \end{itemize}
        以下哪一项陈述代表将资产归类为交易账户资产的标准？
        \begin{itemize}
            \item I.银行必须能够实物交易资产
            \item II.资产的风险必须由银行的交易台管理
        \end{itemize}
    \end{question}
    
    \begin{choices}
        \item I only
        \item II only
        \item Both I and II
        \item Neither I nor II
    \end{choices}
    
    \begin{correctanswer}
        C
    \end{correctanswer}
    
    \begin{explanation}
        \textbf{A选项错误。}
        \begin{itemize}
            \item 仅满足能够实物交易的条件不足以将资产归类为交易账户资产。银行还需要在交易台管理相关风险。
        \end{itemize}
        
        \textbf{B选项错误。}
        \begin{itemize}
            \item 仅满足在交易台管理风险的条件不足以将资产归类为交易账户资产。银行还需要能够对资产进行实物交易。
        \end{itemize}
        
        \textbf{C选项正确。}
        \begin{itemize}
            \item 将资产归类为交易账户资产需要同时满足两个条件：1）银行必须能够对资产进行实物交易；2）银行必须在交易台管理相关风险。这两个条件缺一不可。
        \end{itemize}
        
        \textbf{D选项错误。}
        \begin{itemize}
            \item 这两个条件都是将资产归类为交易账户资产的必要标准，不能同时否定。
        \end{itemize}
    \end{explanation}
    
    \checklist{交易账户资产的分类标准}


    \begin{question}
        Which of the following statements regarding the differences between Basel I, Basel II.5, and the Fundamental Review of the Trading Book (FRTB) for market risk capital calculations is incorrect?
        \\ 以下哪一项陈述关于巴塞尔协议I、巴塞尔协议II.5和交易账簿基本审查（FRTB）在市场风险资本计算方面的差异是不正确的？
    \end{question}
    
    \begin{choices}
        \item Both Basel I and Basel II.5 require calculation of VaR with a 99\% confidence interval\\
        巴塞尔协议I和巴塞尔协议II.5都要求计算99\%置信水平的VaR
        \item FRTB requires the calculation of expected shortfall with a 97.5\% confidence interval\\
        FRTB要求计算97.5\%置信水平的预期短缺
        \item FRTB requires adding a stressed VaR measure to complement the expected shortfall calculation\\
        FRTB要求添加压力VaR度量来补充预期短缺计算
        \item The 10-day time horizon for market risk capital proposed under Basel I incorporates a recent period of time, which typically ranges from one to four years\\
        巴塞尔协议I提出的10天期限的市场风险资本要求包括最近一段时间，通常范围为一年到四年
    \end{choices}
    
    \begin{correctanswer}
        C
    \end{correctanswer}
    
    \begin{explanation}
        \textbf{A选项恰当表述。}
        \begin{itemize}
            \item 巴塞尔协议I和巴塞尔协议II.5都使用99\%置信水平的VaR作为风险度量方法。
        \end{itemize}
        
        \textbf{B选项恰当表述。}
        \begin{itemize}
            \item FRTB使用97.5\%置信水平的预期短缺（Expected Shortfall）作为风险度量方法。
        \end{itemize}
        
        \textbf{C选项不恰当表述。}
        \begin{itemize}
            \item FRTB不需要在预期短缺计算中添加压力VaR。实际上，是巴塞尔协议II.5要求添加压力VaR作为补充风险度量。
        \end{itemize}
        
        \textbf{D选项恰当表述。}
        \begin{itemize}
            \item 巴塞尔协议I中10天期限的市场风险资本要求确实包含了最近一段时间（通常为1-4年）的数据。
        \end{itemize}
    \end{explanation}
    
    \checklist{巴塞尔协议II.5与FRTB在压力VaR要求上的差异}


    \begin{question}
        A large commercial bank is using VaR as its main risk measurement tool. Expected shortfall (ES) is suggested as a better alternative to use during market turmoil. What should be understood regarding VaR and ES before modifying current practices?
        \\ 一家大型商业银行使用VaR作为其主要风险度量工具。预期短缺（ES）被建议作为在市场动荡期间使用的更好替代方法。在修改当前做法之前，应该理解哪些关于VaR和ES？
    \end{question}
    
    \begin{choices}
        \item Despite being more complicated to calculate, ES is easier to backtest than VaR\\
        尽管计算更复杂，但ES比VaR更容易进行回测
        \item Relative to VaR, ES leads to more required economic capital for the same confidence level\\
        与VaR相比，ES在相同置信水平下导致更高的经济资本要求
        \item While VaR ensures that the estimate of portfolio risk is less than or equal to the sum of the risks\\
        虽然VaR确保投资组合风险估计小于或等于各风险之和，但ES考虑了超过置信阈值的损失的严重程度
        \item Both VaR and ES account for the severity of losses beyond the confidence threshold\\
        虽然VaR确保投资组合风险估计小于或等于各风险之和，但ES考虑了超过置信阈值的损失的严重程度
    \end{choices}
    
    \begin{correctanswer}
        B
    \end{correctanswer}
    
    \begin{explanation}
        \textbf{A选项错误。}
        \begin{itemize}
            \item ES实际上比VaR更难进行回测，因为ES需要考虑尾部损失的平均值，而VaR只需要考虑特定分位数。
        \end{itemize}
        
        \textbf{B选项正确。}
        \begin{itemize}
            \item 在相同置信水平下，ES总是大于或等于VaR，因为ES考虑了超过特定置信水平的预期损失的严重程度，而VaR只衡量该置信水平下的最小预期损失。因此，ES会导致更高的经济资本要求。不过，监管机构通常会通过降低使用ES的银行的置信水平要求来调整这种差异。
        \end{itemize}
        
        \textbf{C选项错误。}
        \begin{itemize}
            \item 这个选项不完整，且VaR并不保证投资组合风险估计小于或等于各风险之和。实际上，VaR可能低估尾部风险。
        \end{itemize}
        
        \textbf{D选项错误。}
        \begin{itemize}
            \item 只有ES考虑了超过置信阈值的损失的严重程度，而VaR只关注特定分位数的损失水平。
        \end{itemize}
    \end{explanation}
    
    \checklist{VaR与ES在经济资本要求上的差异}




\section{考点2:Fundamental Review ofthe Trading Book}
\setcounter{question}{0}

\section{考点3:Minimum capital reuuirement for marketrisk}
\setcounter{question}{0}



\end{document}